% !TITLE 行行重行行
% !AUTHOR 无名氏
% !DYNASTY 两汉
% !GENRE 五言古诗
% !ORDER 1

\begin{poem}
行行重行行\textsuperscript{1}，与君生别离。
相去万余里，各在天一涯\textsuperscript{2}。
道路阻且长，会面安可知？
胡马依北风，越鸟巢南枝\textsuperscript{3}。
相去日已远，衣带日已缓\textsuperscript{4}。
浮云蔽白日，游子不顾反\textsuperscript{5}。
思君令人老，岁月忽已晚。
弃捐\textsuperscript{6}勿复道，努力加餐饭。
\end{poem}

\begin{annotations}
\notetext{1}{行行重行行：走了一程又一程。重（chóng），又、再。形容道路遥远无尽，越走越远。}
\notetext{2}{天一涯：天的一角，指极远的地方。涯，边际。}
\notetext{3}{胡马依北风，越鸟巢南枝：北方的马依恋北风，南方的鸟在朝南的树枝上筑巢。这里用动物本能依恋故土，来比喻游子即使身在远方也应思念家乡。}
\notetext{4}{衣带日已缓：衣服的带子一天天变松了。形容因为深深的思念而日渐消瘦。}
\notetext{5}{顾反：回头、回家。反，通“返”。}
\notetext{6}{弃捐：抛弃、丢开。这里指把满腔的愁苦丢开不再提起。}
\end{annotations}

\begin{appreciation}
这是古诗十九首中的第一首，也是离别主题中最动人心魄的篇章之一。

“行行重行行”——走啊走，走个不停。五个字全是动词，像电影里的长镜头，一个人越走越远，没有尽头。“与君生别离”——“生别离”三个字极重。不是普通的分别，而是活着的分离，是生与死的距离。屈原说“悲莫悲兮生别离”，这里直接化用，把离别的痛苦推到了极致。

“相去万余里，各在天一涯”——相隔万里，各自在天边。这是空间的距离。“道路阻且长，会面安可知”——道路又艰险又漫长，什么时候能再见呢？这是时间的距离。空间和时间同时拉远，思念就变成了绝望的等待。

“胡马依北风，越鸟巢南枝”——北方的马依恋北风，南方的鸟在朝南的树枝上筑巢。动物尚且眷恋故土，何况人呢？这里用动物的习性比喻人对故乡和亲人的思念，是中国诗歌中“比兴”手法的经典运用。

“相去日已远，衣带日已缓”——分离的日子越来越长，衣带却越来越松。这不是在减肥，而是因为思念而日渐消瘦。“衣带渐宽终不悔”就是从这儿化出来的。“浮云蔽白日，游子不顾反”——太阳被浮云遮住了，游子不想回来。“白日”比喻光明磊落的心，“浮云”比喻外界的诱惑或阻碍。女子在担心：他是不是在外面遇到了什么，不想回来了？

最后两句是全诗的华彩。“思君令人老，岁月忽已晚”——想你让我变老，时间不知不觉就过去了。一个“忽”字，写出了人在思念中对时间流逝的迟钝。“弃捐勿复道，努力加餐饭”——不要再说了，好好吃饭吧。这看似平淡的祝福，其实藏着最深的爱：不管你回不回来，我都希望你好好的。这种克制的深情，比任何痛哭流涕都更动人。
\end{appreciation}
